\documentclass[12pt,fleqn,a4paper]{report}
%\documentclass[12pt,fleqn,a4paper,draft]{report}
\usepackage[brazil]{babel}
\usepackage[utf8]{inputenc} 
%------------------------------------------------------
\usepackage{caption}
\usepackage{algorithm}
\usepackage{algpseudocode}
%\usepackage[compact]{titlesec}
\usepackage{indentfirst}
%------------------------------------------------------

\usepackage{graphicx}
\usepackage{amssymb}
\usepackage{amsmath,amsfonts,amscd,bezier,amstext}
\usepackage{colortbl}
\usepackage{setspace}
\usepackage{pxfonts}
\usepackage{subfigure,float,subfloat}
\usepackage{textcomp}
\usepackage{boxedminipage}
\usepackage{algorithm}
\usepackage[brazil,pagebackref,colorlinks,linkcolor=black,citecolor=black]{hyperref}
\usepackage{multirow}
\usepackage{booktabs}
\usepackage{textcomp}
\usepackage{enumitem}
\usepackage{mathtools}
\usepackage[title]{appendix}
\usepackage{animate}
\usepackage{rotating} %% permite rotacionar objetos do tipo ``floating''
%\usepackage{subfig}
\usepackage[flushleft]{threeparttable}
\usepackage{xcolor}
\usepackage{microtype}
\usepackage{icomma}
\usepackage{fancyvrb}
\usepackage{fancyhdr}
\usepackage{longtable}
\usepackage{caption}
\usepackage{stackengine}
\usepackage{nomencl}
\usepackage{ifthen}
\usepackage{makecell}
\usepackage{scalerel}
\usepackage{here}
\usepackage{marvosym}			% Alguns símbolos
\usepackage{wasysym}			% Alguns símbolos
\usepackage[alf,abnt-etal-list=0,abnt-etal-text=it,abnt-emphasize=it,bibjustif,abnt-last-names=bibtex,abnt-url-package=url]{abntex2cite}
\usepackage{filecontents}
\usepackage{lscape}
\usepackage{color}
\usepackage{pgf}
\usepackage{tikz}
\usepackage{ragged2e}
%\usepackage{cleveref}
%\usepackage[brazilian,hyperpageref]{backref}	 % Paginas com as citações na bibl

% ---
% Configurações do pacote backref
% Usado sem a opção hyperpageref de backref
\renewcommand{\backrefpagesname}{Citado na(s) página(s):~}
% Texto padrão antes do número das páginas
\renewcommand{\backref}{}
% Define os textos da citação
\renewcommand*{\backrefalt}[4]{
	\ifcase #1 %
		Nenhuma citação no texto.%
	\or
		Citado na página #2.%
	\else
		Citado #1 vezes, nas páginas #2.%
	\fi}%
% ---

\usetikzlibrary{shapes,arrows}
%
%\pdfminorversion=5 
%\pdfcompresslevel=9
%\pdfobjcompresslevel=9
%\setlength{\textheight}{24cm} \setlength{\textwidth}{16cm}
%\voffset=-1.5cm \hoffset=-1.2cm
\DeclareGraphicsExtensions{.JPG,.jpg,.tiff,.pdf,.mps,.PNG,.png}

%\newfloat{algorithm}{!ht}{loa}
%\floatname{algorithm}{Algorithm}

\DeclareMathOperator{\sen}{sen}
\usepackage[normalem]{ulem}

\usepackage{color}
\definecolor{dkgreen}{rgb}{0,0.6,0}
\definecolor{gray}{rgb}{0.5,0.5,0.5}
\definecolor{mauve}{rgb}{0.58,0,0.82}
\newcommand{\mltc}[3]{\multicolumn{#1}{#2}{#3}}
\newcommand{\mltr}[2]{\multirow{#1}{*}{#2}}
\newcommand{\done}[2]{\frac{\partial{#1}}{\partial{#2}}}
\newcommand{\dtwo}[2]{\frac{\partial{^2#1}}{\partial{#2}^2}}
\newcommand{\dtwotwo}[3]{\frac{\partial{^2#1}}{\partial{#2}\partial{#3}}}
\newcommand{\dn}[3]{\frac{\partial{^{#3}#1}}{\partial{#2}^{#3}}}

%=============================================================
%============== Formatação dos códigos =======================
%=============================================================
% ---> Inserir códigos
\usepackage{listings}		
% ---> Aceitar acentos nos códigos
\usepackage{listingsutf8}		 
\lstset{inputencoding=utf8/latin1}

% ---> Definição dos acentos para o listings:
\lstset{
	literate=%
	{é}{{\'{e}}}1
	{è}{{\`{e}}}1
	{ê}{{\^{e}}}1
	{ë}{{\¨{e}}}1
	{É}{{\'{E}}}1
	{Ê}{{\^{E}}}1
	{û}{{\^{u}}}1
	{ù}{{\`{u}}}1
	{ú}{{\'{u}}}1
	{â}{{\^{a}}}1
	{à}{{\`{a}}}1
	{á}{{\'{a}}}1
	{ã}{{\~{a}}}1
	{Á}{{\'{A}}}1
	{Â}{{\^{A}}}1
	{Ã}{{\~{A}}}1
	{ç}{{\c{c}}}1
	{Ç}{{\c{C}}}1
	{õ}{{\~{o}}}1
	{ó}{{\'{o}}}1
	{ô}{{\^{o}}}1
	{Õ}{{\~{O}}}1
	{Ó}{{\'{O}}}1
	{Ô}{{\^{O}}}1
	{î}{{\^{i}}}1
	{Î}{{\^{I}}}1
	{í}{{\'{i}}}1
	{Í}{{\~{Í}}}1
}

\lstset{language=[90]Fortran,
	basicstyle=\tiny,
	numbers=left,
	frame=shadowbox,
	showstringspaces=false,
	rulesepcolor=\color{black}}

\renewcommand{\lstlistingname}{Código}

%=============================================================
%=============================================================

\pagestyle{fancy}
\fancyhf{}
\usepackage[top=2.0 cm,left=2 cm,right=2 cm,bottom=2.0 cm,includeheadfoot,headheight=30pt]{geometry} %Formata as margens do textos
\rhead{\includegraphics[width=1cm]{layout_document/ufu.pdf}}
\lhead{Relatório Técnico Parcial 2}
\rfoot{ \thepage}
% Redefine the plain page style
\fancypagestyle{plain}{
	\fancyhf{}
	\rfoot{ \thepage}
	\rhead{\includegraphics[width=1cm]{layout_document/ufu.pdf}}
	\lhead{Relatório Técnico Parcial 2}
}
\fancypagestyle{capa}{
	\fancyhf{}
}

% Length to control the \fancyheadoffset and the calculation of \headline
% simultaneously
\newlength\FHoffset
\setlength\FHoffset{0cm}

\addtolength\headwidth{2\FHoffset}

\fancyheadoffset{\FHoffset}

% these lengths will control the headrule trimming to the left and right 
\newlength\FHleft
\newlength\FHright

% here the trimmings are controlled by the user
\setlength\FHleft{0cm}
\setlength\FHright{0cm}

%\pagestyle{fancy}
%\fancyhf{}
%%\usepackage[top=2.0 cm,left=3.5 cm,right=3.0 cm,bottom=2.0 cm,includeheadfoot]{geometry} %Formata as margens do textos
%\usepackage[top=2.0 cm,left=2 cm,right=2 cm,bottom=2.0 cm,includeheadfoot,headheight=30pt]{geometry} %Formata as margens do textos
%\rhead{\includegraphics[width=1cm]{layout_document/ufu.pdf}}
%\lhead{Relatório técnico parcial}
%\rfoot{ \thepage}
%% Redefine the plain page style
%\fancypagestyle{plain}{
%	\fancyhf{}
%	\rfoot{ \thepage}
%	\rhead{\includegraphics[width=1cm]{layout_document/ufu.pdf}}
%	\lhead{Relatório técnico parcial}
%}
%\fancypagestyle{capa}{
%	\fancyhf{}
%}

\addto\captionsportugues{% Replace "english" with the language you use
	\renewcommand{\contentsname}%
	{Sumário}%
}

\newcommand{\figuraabnt}[6]{
	\begin{figure}[H]
		\centering
		\caption{#1}
		\centering
		\begin{tabular}{c}
			{\includegraphics[trim={#2},clip,width=#3\linewidth]{#4}}\\
			\multicolumn{1}{l}{\footnotesize{Fonte: #5}}\\
		\end{tabular}
		\label{#6}
	\end{figure}
}

\newcommand{\duasfigurasabnt}[9]{
	\begin{figure}[H]
		\centering
		\caption{#1}
		\centering
		\begin{tabular}{c c}
			{\includegraphics[width= #3 \linewidth]{#4}} & {\includegraphics[width= #6 \linewidth]{#7}}\\
			{#2}&{#5}\\
			\multicolumn{2}{l}{\footnotesize{Fonte: #8}} \\
		\end{tabular}
%	    \legend{Fonte: #8}
		\label{#9}
	\end{figure}
	
}


\usepackage{float}
\usepackage{caption}

\newfloat{grafico}{htpb}{grf}[figure]
\floatstyle{plaintop}
\restylefloat{grafico}
\renewcommand\thegrafico{\arabic{grafico}}
\newcommand{\graficoname}{Gr\'{a}fico}
\floatname{grafico}{\graficoname}

\newcommand{\listofgraficoname}{Lista de gráficos}
\newcommand{\listofgrafico}{%
%	\addcontentsline{toc}{chapter}{\listofgraficoname}
	\listof{grafico}{\listofgraficoname}
}
\newcounter{graficocount}

\newcommand{\doisgraficosabnt}[9]{
	
	 \setcounter{grafico}{\value{graficocount}}
	 
	\begin{grafico}[H]

		\caption[Gráfico \ref{#9}~--~#1]{#1\label{#9}}
		\centering
		\begin{tabular}{c c}
			{\includegraphics[width= #3 \linewidth]{#4}} & {\includegraphics[width= #6 \linewidth]{#7}}\\
			{\footnotesize{#2}}&{\footnotesize{#5}}\\
			\multicolumn{2}{l}{\footnotesize{Fonte: #8}}
		\end{tabular}
		
	\end{grafico}
	\stepcounter{graficocount}
}


\newcommand{\doisgraficosabntcoluna}[9]{
	
	\setcounter{grafico}{\value{graficocount}}
	
	\begin{grafico}[H]
		
		\caption[Gráfico \ref{#9}~--~#1]{#1\label{#9}}
		\centering
		\begin{tabular}{c}
			{\includegraphics[width= #3 \linewidth]{#4}}\\
			{\footnotesize{#2}}\\
			{\includegraphics[width= #6 \linewidth]{#7}}\\
			{\footnotesize{#5}}\\
			\multicolumn{1}{l}{\footnotesize{Fonte: #8}}
		\end{tabular}
		
	\end{grafico}
	\stepcounter{graficocount}
}



\newcommand{\graficoabnt}[5]{
	 \setcounter{grafico}{\value{graficocount}}
	\begin{grafico}[H]
		
		\caption[Gráfico \ref{#5}~--~#1]{#1\label{#5}}
		\centering
		\begin{tabular}{c c}
			\multicolumn{2}{c}{\includegraphics[width= #2 \linewidth]{#3}}\\
			\multicolumn{2}{l}{\footnotesize{Fonte: #4}}
		\end{tabular}
		
	\end{grafico}
	\stepcounter{graficocount}
}
\newcommand{\nasa}{elaborada pelo autor.}
\newcommand{\nasas}{elaboradas pelo autor.}
\newcommand{\naso}{elaborado pelo autor.}
\newcommand{\nasos}{elaborados pelo autor.}

\newcommand{\nasavisit}{\nasa{} Captura de tela do software VisIt \cite{visit}.}
\newcommand{\nasasvisit}{\nasas{} Capturas de tela do software VisIt \cite{visit}.}
\newcommand{\nasaanflex}{\nasa{} Captura de tela do software Anflex \cite{anflex}.}

\newcommand{\nasapara}{\nasa{} Captura de tela do software Paraview \cite{paraview}.}
\newcommand{\nasaspara}{\nasas{} Capturas de tela do software Paraview \cite{paraview}.}

\newcommand{\nasomatplot}{\naso{} Gráfico gerado com Matplotlib \cite{matplotlib}.}
\newcommand{\nasosmatplot}{\nasos{} Gráficos gerados com Matplotlib \cite{matplotlib}.}

\newcommand{\nasamatplot}{\nasa{} Figura gerada com Matplotlib \cite{matplotlib}.}
\newcommand{\nasasmatplot}{\nasas{} Figuras geradas com Matplotlib \cite{matplotlib}.}

\newcommand{\Fonte}[2]{\multicolumn{#1}{l}{\footnotesize{Fonte: #2}}}
\newcommand{\Fontefig}[1]{\flushleft\footnotesize{Fonte: #1}}
\newcommand{\adaptado}[1]{adaptado de \citeonline{#1}.}
\newcommand{\adaptada}[1]{adaptada de \citeonline{#1}.}

%%%%%%%%%%%%%%%%FIGURAS%%%%%%%%%%%%%%

%Ao longo do texto

\newcommand{\fig}[1]{\hyperref[{#1}]{Fig. \ref{#1}}}
\newcommand{\figs}[2]{Figs. \ref{#1} e \ref{#2}}
\newcommand{\figstres}[3]{Figs. \ref{#1}, \ref{#2} e \ref{#3}}
\newcommand{\figsquatro}[4]{Figs. \ref{#1}, \ref{#2}, \ref{#3} e \ref{#4}}

%Início de frase:

\newcommand{\figini}[1]{\hyperref[{#1}]{Figura \ref{#1}}}
\newcommand{\figsini}[2]{Figuras \ref{#1} e \ref{#2}}
\newcommand{\figsiniquatro}[4]{Figuras \ref{#1}, \ref{#2}, \ref{#3} e \ref{#4}}

%%%%%%%%%%%%%%%%GRAFICOS%%%%%%%%%%%%%%

%Ao longo do texto

\newcommand{\graf}[1]{\hyperref[{#1}]{Gráf. \ref{#1}}}
\newcommand{\grafs}[2]{Gráfs. \ref{#1} e \ref{#2}}

%Início de frase:

\newcommand{\grafini}[1]{\hyperref[{#1}]{Gráfico \ref{#1}}}
\newcommand{\grafsini}[2]{Gráficos \ref{#1} e \ref{#2}}


%%%%%%%%%%%%%%%%TABELAS%%%%%%%%%%%%%%

%Ao longo do texto

\newcommand{\tab}[1]{\hyperref[{#1}]{Tab. \ref{#1}}}
\newcommand{\tabs}[2]{Tabs. \ref{#1} e \ref{#2}}

%Início de frase:

\newcommand{\tabini}[1]{\hyperref[{#1}]{Tabela \ref{#1}}}
\newcommand{\tabsini}[2]{Tabelas \ref{#1} e \ref{#2}}

%%%%%%%%%%%%%%%%EQUAÇÕES%%%%%%%%%%%%%%

%Ao longo do texto

\newcommand{\eq}[1]{\hyperref[{#1}]{Eq. (\ref{#1})}}
\newcommand{\eqsa}[2]{Eqs. (\ref{#1}) a (\ref{#2})}
\newcommand{\eqs}[2]{Eqs. (\ref{#1}) e (\ref{#2})}
\newcommand{\eqsquatro}[4]{Eqs. (\ref{#1}), (\ref{#2}), (\ref{#3}) e (\ref{#4})}
\newcommand{\eqsoito}[8]{Eqs. (\ref{#1}), (\ref{#2}), (\ref{#3}), (\ref{#4}), (\ref{#5}), (\ref{#6}), (\ref{#7}), (\ref{#8})}

%Início de frase:

\newcommand{\eqini}[1]{\hyperref[{#1}]{Equação (\ref{#1})}}
\newcommand{\eqsini}[2]{Equações (\ref{#1}) e (\ref{#2})}

%%%%%%%%%%%%%%%%ANEXO%%%%%%%%%%%%%%

%\newcommand{\refanexo}[1]{ANEXO \ref{#1}}

\newcommand*{\refanexo}[1]{\hyperref[{#1}]{ANEXO~\ref*{#1}}}
%%%%%%%%%%%%%%%%APÊNDICE%%%%%%%%%%%%%%

%\newcommand{\refapendice}[1]{AP\^{E}NDICE \ref{#1}}

\newcommand*{\refapendice}[1]{\hyperref[{#1}]{AP\^{E}NDICE~\ref*{#1}}}

%%%%%%%%%%%%%%%%Seção%%%%%%%%%%%%%%

\newcommand*{\refsec}[1]{\hyperref[{#1}]{seção~\ref*{#1}}}
\newcolumntype{C}[1]{>{\centering\let\newline\\\arraybackslash\hspace{0pt}}m{#1}}

\newcommand{\stkout}[1]{\ifmmode\text{\sout{\ensuremath{#1}}}\else\sout{#1}\fi} % cortar uma letra

\newcommand{\reg}[1]{#1$^{\tiny{\circledR}}$} % Símbolo de marca registrada

\newcommand{\MFSim}{MFSim}

\newcommand{\PETSc}{\textit{PETSc}}
%\include{new_commands/zoom_in_figure}

\begin{document}

	\baselineskip 25pt
	\thispagestyle{empty}
	
	\begin{singlespace}
	
		\begin{figure}[htb]
			\begin{center}
				\includegraphics[scale=0.7]{layout_document/ufu.pdf}\\
				{\bf Universidade Federal de Uberl\^andia}
			\end{center}
		\end{figure}
		
		\begin{figure}[htb]
			\begin{minipage}{0.49\linewidth}
				\centering
				\includegraphics[scale=0.15]{layout_document/lmest.pdf}\\
				{\bf \small{Laborat\'orio de Mec\^anica de Estruturas}} \\
				{\bf \small{Jos\'e Eduardo Tann\'us Reis}}
			\end{minipage}
			\begin{minipage}{0.49\linewidth}
				\centering
				\vspace{-0.15cm}\includegraphics[scale=0.22]{layout_document/mflab.pdf}\\ \vspace{0.15cm}
				{\bf \small{Laborat\'orio de Mec\^anica dos Fluidos}}
			\end{minipage}
		\end{figure}
	
	\end{singlespace}
	%\vspace*{30 pt}
	
	\begin{large}
		\begin{center}
			Relatório Técnico Parcial 2\\
			\vspace*{10 pt}
			Período: 01/01/2020 à 31/08/2020\\
			\vspace*{25 pt}
			{\bf Desenvolvimento de uma ferramenta computacional dedicada à construção de metamodelos de estruturas submersas}
		\end{center}
	\end{large}
	
	\vspace*{10 pt}
	\begin{normalsize}
		\begin{description}
			\item[] 
	
			\hspace{0.7cm} UNIVERSIDADE FEDERAL DE UBERLÂNDIA - UFU\\
			FUNDA\c C\~AO DE APOIO UNIVERSITÁRIO - FAU
			\begin{flushleft}
				Termo de Cooperação: PT-127.01.13870 \\
				Número SAP: 4600585313 \\
				Número SIC: 2019/00056-2 \\
				Número jurídico: 5900.0111467.19.9\\
				Data de Assinatura: 31 de maio de 2019\\
				Data de Término Prevista: 29 de maio de 2021
			\end{flushleft}
	
			\vspace{10pt}
			\item[Coordenador:] Prof. Dr. Aldemir Aparecido Cavalini Jr  
			\vspace{10pt}
			\item[Gerencia Técnica:] CENPES/PDDP/TDUT
			\item[Interlocutor Técnico:] Marcos André Duarte Martins
	
		\end{description}
	\end{normalsize}
	
	\newpage
	
	\section*{Equipe Responsável}

		\begin{normalsize}
			\begin{description}
				\vspace{20pt}
				\item[Coordenador:] Prof. Dr. Aldemir Aparecido Cavalini Jr 
				\item[\vspace{10pt}]  
		
				\item[Pesquisadores:] Prof. Dr. Aristeu da Silveira Neto \\
				\hspace*{1.95cm} Prof. Dr. João Marcelo Vedovoto \\
				\hspace*{1.95cm} Dra. Denise Kinoshita \\
				\hspace*{1.95cm} Dr. Marcelo Maia Ribeiro Damasceno \\
				\hspace*{1.95cm} Msc. Hélio Ribeiro Neto \\
				\hspace*{1.95cm} Msc. Tatiane Nunes da Costa \\
				\hspace*{1.95cm} Eng. André Rezende Dessimoni Carvalho \\
				\hspace*{1.95cm} Eng. Bruno Resende Ferreira Rende\\
				\hspace*{1.95cm} Alessandra Lopes Paiva\\
				\hspace*{1.95cm} Arturo Burgos Moreira\\
				\hspace*{1.95cm} Felipe José Oliveira Ribeiro\\
		
				\item[Responsável da Convenente:] Gisele Debortoli Guilherme
				 
				\item[Gerencia Técnica:] CENPES/PDDP/TDUT
				\item[Interlocutor Técnico:] Marcos André Duarte Martins
			\end{description}
		\end{normalsize}

		\newpage
	\section*{Resumo}
	
		\noindent O presente relatório técnico parcial corresponde às atividades realizadas durante o período de janeiro de 2020 -- agosto de 2020, a respeito do projeto \textit{Desenvolvimento de uma Ferramenta Computacional Dedicada à Construção de Metamodelos de Estruturas Submersas} sob o termo de cooperação PT-127.01.13870 / SAP 4600585313. Devido ao caráter multidisciplinar do presente projeto de pesquisa, este relatório é dividido em capítulos (quase) independentes, cada qual com sua metodologia, discussões e resultados específicos. Como regra geral, os desenvolvimentos seguiram as linhas: \textcolor{red}{\textbf{CONCLUIR CONCLUIR CONCLUIR CONCLUIR CONCLUIR CONCLUIR CONCLUIR}}.
	
	\newpage
	\tableofcontents
	\newpage
	
	%\include{informacoes}
	%\include{objetivos}
	
	%\setcounter{page}{2}
	%\tableofcontents \pagestyle{plain}
	
	\chapter{Introdução}\label{chap_intro}

%\pagenumbering{arabic} % numerar as paginas em algarismos arabicos (ou romanos)
\thispagestyle{empty} %Oculta o numero da primeira pagina do capitulo
\vspace{3ex}

A dinâmica dos fluidos e das estruturas imersas são ramos da Engenharia Mecânica com as quais se estuda o comportamento das interações fluido-estruturais, visando aperfeiçoar vários processos na indústria, assim como entender fenômenos na natureza. O estudo desse tipo de fenômeno pode ser feito de duas maneiras: a experimentação material e a experimentação virtual. Em ambas as frentes é necessária a modelagem física do problema em análise e a montagem de bancadas experimentais que representem a física do problema a ser estudado. Para a modelagem material, faz-se necessária a construção de bancadas e as respectivas instrumentações para permitir a coleta de dados e análise estatística dos resultados obtidos. A modelagem física consiste na avaliação do problema de interesse e determinação de suposições físicas que visam viabilizar a análise. Na segunda maneira, é necessária a modelagem física e matemática do problema de interesse. A modelagem matemática consiste na obtenção de equações diferenciais, integrais e/ou integro-diferenciais que modelam a física assciada e depois utilizar métodos numéricos apropriados para a discretização das equações. 

Os métodos de análise computacional estão em franco crescimento. É importante destacar que o método computacional não substitui o experimental material, mas o complementa. Além disso, esta metodologia é muito versátil, permitindo uma análise minuciosa do problema físico e uma maior flexibilidade em relação às condições físicas. Alguns experimentos podem ser perigosos de se reproduzir em laboratório, ou até mesmo impossíveis. Algumas desvantagens são a necessidade de modelos numéricos adequados e de computadores robustos, conforme o problema que se está analisando.

A interação entre escoamentos e estruturas é um problema complexo e recorrente em aplicações de engenharia. Esse fenômeno pode ser encontrado em aeronaves, motores a jato, dutos, reatores nucleares e químicos, pontes, torres, plataformas \textit{off-shore}, válvulas de compressores, coração, aneurismas, entre outros. A dinâmica dos fluidos computacional, aliada à solução numérica das equações que modelam a movimentação de estruturas, é uma grande aliada na compreensão dos problemas de interação fluido-estrutura. Trata-se de um problema multidisciplinar, visto que envolve, por exemplo, a mecânica dos fluidos, mecânica das estruturas, engenharia de software e a ciência da computação.

Os escoamentos sobre estruturas cilíndricas podem ser a fonte de vibrações induzidas por estruturas turbilhonares. Essas vibrações podem induzir um aumento das forças fluidodinâmicas, ou seja, arrasto e sustentação, levando assim a um aumento dos esforços aplicados sobre as estruturas. Alguns resultados indicam que escoamentos bidimensionais sobre cilindros circulares mudam, por exemplo, o coeficiente de arrasto médio de 1,3 para 2,2. Além disso, as vibrações podem causar nucleação e propagação de trincas na estrutura conduzindo-a a falha em virtude da fadiga. Em alguns casos, o valor RMS do coeficiente de sustentação pode ser alterado de 0,3 para 1,75, dependendo do regime de operação \cite{Chern2014}. Esses são resultados que justificam a preocupação com o processo de interação fluido-estrutura em cilindros. Isso é especialmente importante quando esses cilindros são dutos pelos quais petróleo ou gás natural são transportados, sobre os quais se têm ondas e/ou correntes marítimas atuando. A manutenção desse tipo de duto é cara, visto que podem estar a centenas de metros da superfície. Qualquer falha nessas estruturas pode causar desastres ambientais e grandes prejuízos. Por isso é importante entender como o processo de interação fluido-estrutura atua sobre a dinâmica do duto para prevenir falhas.

A pesquisa apresentada neste projeto é resultado de cooperação entre o Laboratório de Mecânica dos Fluidos (MFLab), o Laboratório de Mecânica de Estruturas Prof. José Eduardo Tannús Reis (LMEst) da Universidade Federal de Uberlândia (UFU) e com o centro de pesquisa (CENPES) da Petróleo Brasileiro S.A. (Petrobras). A pesquisa está sendo feita utilizando a ferramenta computacional MFSim, que está sendo desenvolvido no MFLab. Nessa ferramenta é possível a execução de simulações de escoamentos incompressíveis levando em consideração a movimentação de estruturas. Para isso, o método da Fronteira Imersa é utilizado. Esta metodologia é particularmente adequada para os problemas que envolvem interação fluido-estrutura, pois permite tratar os domínios do fluido e da estrutura de forma independente. As equações que modelam o escoamento são resolvidas em um domínio euleriano fixo e cartesiano, enquanto a superfície do corpo imerso (estrutura) é representada por um conjunto de pontos lagrangeanos. Através dessa técnica, as forças na interface entre a estrutura e o fluido são avaliadas e utilizadas tanto nos códigos associados ao fluido para imposição da condição de contorno de não deslizamento, quanto na rotina estrutural para o cálculo dos deslocamentos e velocidades da estrutura.

A motivação para a realização deste projeto de pesquisa se dá pelo fato que é alto o custo computacional de simulações associadas ao dimensionamento de estruturas imersas. Este problema se torna maior caso a aplicação de processos de otimização seja necessária. Assim, neste projeto é desenvolvida uma ferramenta computacional dedicada à construção de metamodelos dessas estruturas. Ao final do projeto, a análise dinâmica de diferentes estruturas imersas poderá ser realizada com baixo custo computacional. Neste caso, em vez de resolver as equações do movimento associadas ao problema de interação fluido-estrutura, os usuários da ferramenta acessarão um mapa criado com funções e polinômios obtidos a partir de simulações de alto custo realizadas previamente. Desta forma, reduzindo o tempo computacional para alguns segundos de cálculo viabilizando a utilização de mais testes e avaliações no projeto estrutural. 

Neste projeto, serão realizadas atividades voltadas para o desenvolvimento de uma ferramenta computacional usada na construção de metamodelos para análise estrutural. Assim, diferentes análises e procedimentos de otimização poderão ser realizados de forma mais rápida e eficiente no projeto deste tipo de estruturas. Além disso, será realizada uma análise computacional de escoamentos turbulentos tridimensionais sobre um duto com \textit{strakes}. Desta forma, ao final do projeto, a análise de diferentes estruturas submersas será realizada com baixo custo computacional. Neste sentido, será utilizada uma plataforma desenvolvida no MFlab que vem sendo empregada em aplicações de interesse da indústria de óleo e gás. Novos desenvolvimentos nesse código serão necessários. Com esses novos desenvolvimentos a ferramenta será potencializada para a aplicação em pauta. Esta é a importância deste projeto para o setor de petróleo, gás natural, energia e biocombustíveis.

A presente proposta tem por objetivo dar sequência à cooperação com a construção da ferramenta de metamodelagem para estruturas de \textit{risers}. Além disso, será realizada a simulação numérica de escoamentos tridimensionais sobre dutos com \textit{strakes}. As dimensões dos \textit{strakes}, ou seja, a altura e passo da hélice, são parâmetros críticos no projeto destes sistemas. Assim sendo, torna-se interessante avaliar sua eficiência através de simulações computacionais para diferentes condições de operação. São objetivos específicos: 1) Metamodelagem estrutural: a parametrização do metamodelo é dada pelo conjunto de dados estruturais e hidrodinâmicos, mas a metamodelagem constitui apenas a análise estrutural. O carregamento hidrodinâmico é dado de entrada; 2) Geração de sinais de CD e CL através de CFD pelo MFSim sobre os modelos que temos de cilindros fixos com strakes, ensaiados em canal de água corrente: análises de CFD puras; 3) Metamodelagem para levantamento de cargas sobre cilindros com \textit{strakes} sem acoplamento estrutural (metamodelagem apenas de CFD).

Neste contexto, o presente relatório parcial é organizado da seguinte maneira. \textcolor{red}{\textbf{CONCLUIR CONCLUIR CONCLUIR CONCLUIR CONCLUIR CONCLUIR CONCLUIR}}.
%	\include{andre/}
%	\include{arturo/}
%   \include{denise/}
	\thispagestyle{empty} %Oculta o numero da primeira pagina do capitulo
\vspace{3ex}


\chapter{Estudo de viabilidade do método KPLS na predição de vida em fadiga.}

\section{Introdução}
% O método KPLS (Kriging e médios quadrados parciais), quando usado na predição de séries temporais de tensão, pode resultar em boas aproximações a baixos custos computacionais.

Neste segmento estuda-se os impactos da utilização do método KPLS (\emph{Kriging model and partial least squares}) na predição de séries temporais de tensão para aplicação no estudo da vida em fadiga dos risers. No caso deste estudo, a análise é feita a partir das tensões externas de Von Mises desenvolvidas a partir do software Anflex do CENPES/PETROBRAS. O caso analisado consiste em um longo \emph{riser} com flutuadores (Fig. \ref{figure.1}) que tornam flutuante parte da estrutura submersa (configuração \emph{lazy-wave}). 

\begin{figure}[!ht]
    \centering
    \includegraphics[width=0.5\linewidth , trim = {40mm 0mm 30mm 5mm} , clip]{felipe/fig_felipe/caso_estudado.jpg}
    \caption{Configuração \emph{lazy-wave}, onde o estudo foi desenvolvido.}
    \label{figure.1}
\end{figure}

O segmento flutuante (denominado região de SAG/HOG) tem o objetivo de atenuar os carregamentos advindos dos deslocamentos da estrutura flutuante no segmento do \emph{riser} em contato com o leito marinho. Seus efeitos podem ser observados no andamento deste estudo ao causar notável influencia na vida em fadiga dos elementos próximos. 

A biblioteca SMT foi utilizada para aplicação do método KPLS. Este pacote possui muitas ferramentas de metamodelo também aplicadas neste projeto \cite{BOUHLEL2019}.

Dessa forma, tem-se o intúito de se aplicar o método \emph{RainFlow} e o método de Goodman modificado para determinação da vida em fadiga a partir das séries de tensões externas de Von Mises advindas das simulações feitas no Anflex. O método KPLS é utilizado para predição parcial destas séries e então se calcula a vida em fadiga destas séries parcialmente ajustadas. Ao comparar-se os resultados observou-se boa conformidade, o que levanta boas oportunidades de otimização ao tornar facultativo o desenvolvimento de uma simulação completa para este tipo de análise. 

\section{Metodologia}

Para o estudo da vida em fadiga se utilizou o método \emph{RainFlow} e o método de Goodman modificado, enquanto que o método utilizado para se prever as séries temporais foi o \emph{KPLS}. 

\subsection{Método RainFlow}\label{rainflow}
Desenvolvido em 1968 por Tatsuo Endo e M. Matsuishi segundo \cite{dowling2007}, o método \emph{Rainflow} tem o objetivo de identificar os ciclos individuais de carregamento. Estas informações, ao se analisar séries temporais periódicas, podem ser usadas no desenvolvimento da vida em fadiga da estrutura analisada, isto é, o número de ciclos que o objeto de estudo aguentará antes de uma fratura devido à fadiga ocorrer. Ele foi desenvolvido com inspiração nos \emph{pagodes}, que são construções de arquitetura asiática muito utilizados como templos budistas (Fig.\ref{pagode}) em países como o Japão, China e as Coreias. 

\begin{figure}[!ht]
    \centering
    \includegraphics[width=0.25\linewidth , trim = {0mm 0mm 0mm 0mm} , clip]{felipe/fig_felipe/pagode.jpg}
    \caption{Pagode asiático.}
    \label{pagode}
\end{figure}

O método se baseia na analogia de um dia de chuva sobre um \emph{pagode}. Ao deixar um dos telhados a água cai verticalmente, podendo ou não tocar os telhados subsequentes. Tocar ou não o andar de baixo tem carácter classificatório sobre um ponto de vista de ciclos de carregamento. Ao se analisar três pontos em uma série temporal (pico, vale e pico) e se imaginar a situação descrita pela analogia, se a água que deixa o primeiro pico não tocar o segundo não há um ciclo de carregamento. Se, por outro lado, tocar, tem-se então identificado um ciclo de carregamento (Fig.\ref{rain_flow_1}). O mesmo pensamento pode ser aplicado em situações de vale, pico e vale.

\begin{figure}[!ht]
    \centering
    \includegraphics[width=0.4\linewidth, angle=90 , trim = {0mm 0mm 0mm 0mm} , clip]{felipe/fig_felipe/rainflow_1.png}
    \caption{Analogia do \emph{pagode} aplicada em um exemplo matemático.}
    \label{rain_flow_1}
\end{figure}

Dessa forma é oportuno pontuar o método com rigor matemático. A análise \emph{Rainflow} é aplicada a conjuntos de três pontos ($X$ , $Y$ e $Z$). Assim, quando $\left|\Delta_{YZ} \right|$ é maior ou igual a $\left|\Delta_{XY} \right|$, tem-se um ciclo de carregamento (Fig.\ref{rain_flow_2}): 

\begin{figure}[!ht]
    \centering
    \includegraphics[width=0.7\linewidth, trim = {0mm 0mm 0mm 0mm} , clip]{felipe/fig_felipe/rain_flow_2.jpg}
    \caption{Definição matemática formal do método de identificação.}
    \label{rain_flow_2}
\end{figure}

Como o método só lida com ciclos de tensão, o primeiro passo nesta análise é excluir da série todos os pontos que não são picos nem vales. Como se assume que este ciclo de carregamento será aplicado de forma contínua e periódica, não faz diferença o ponto de início da análise se os pontos anteriores a este forem transpostos para o final da série temporal de tensões. Assim, inicia-se o método no máximo valor de tensão da série. Isso é feito para evitar problemas com o método de contagem. Além disto repete-se o primeiro ponto (máximo global) no final da série, também para fins de tornar a série amigável ao algorítmo de contagem, sem comprometimentos quanto à acurácia da contagem.

Após a adequação, o algorítimo passará pelos seguintes passos a partir do primeiro elemento da série:
\begin{itemize}
    \item É executada a análise \emph{Rainflow} de três pontos ($X_i$ , $X_{i+1}$ e $X_{i+2}$);
    \item Se identificado um ciclo de carregamento, são armazenados os dados do ciclo ( $\sigma_{m}$ , $\sigma_{a}$ ) e se exclui da análise os dois primeiro pontos ($X_i$ e $X_{i+1}$). Depois disso, o algorítmo é encaminhado para o primeiro elemento da série temporal;
    \item Se não for identificado um ciclo de carregamento a análise é levada ao próximo passo de tempo e se repete este protocolo;    
\end{itemize}

Dessa forma este protocolo se repete até o momento em que sobrar apenas dois pontos, o inicial e o final, momento em que o algorítmo se encerra.
 
As informações que se tira de cada ciclo de tensão identificado são $\sigma_{a}$ e $\sigma_{m}$ que podem ser definidos como:

\begin{equation}\label{eq.1}
    \sigma_a = \left| \frac{\sigma_{i} - \sigma_{i+2}}{2} \right|,
\end{equation}

\begin{equation}\label{eq.2}
    \sigma_m = \frac{\sigma_{i} + \sigma_{i+2}}{2}.
\end{equation}

\subsection{Método de Goodman modificado}

Com o método de Goodman modificado somos capazes de determinar uma tensão alternada equivalente $\sigma_{ar}$, onde o efeito de tensão média $\sigma_{m}$ é levado em conta. Assim, a partir disto determina-se a curva de tensão versus número de ciclos (S-N) que pode ser utilizada para o cálculo da vida em fadiga.

Assim, $\sigma_{ar}$ pode ser comparado diretamente com $\sigma_{a}$ (calculado na Eq.\ref{eq.2}) ao se considerar $R = -1$. Pode-se determinar $\sigma_{ar}$ com o método de Goodman modificado, como descrito em \cite{dowling2007}:

\begin{equation}
    \sigma_{ar} = \frac{\sigma_a}{1 - \frac{\sigma_m}{\sigma_f^\prime}},
\end{equation}

onde $\sigma_f^\prime$ é uma tensão ajustada com base na curva S-N para $\sigma_m = 0$, sendo específica para cada material e constante. Para se chegar a esse valor utiliza-se da seguinte relação:

\begin{equation}
    \sigma_a = \sigma_f^\prime \left( 2 N_f \right)^{b},
\end{equation}

onde $b$ também é encontrado a partir da curva S-N e é única e constante para cada geometria e $N_f$ denota o número de ciclos para a falha do material.

Dessa forma, para determinar este valor, por fim, utiliza-se a seguinte expressão:


\begin{equation}
    N_f = \frac{1}{2} \left( \frac{\sigma_{ar}}{\sigma_f^\prime}  \right)^{1/b}.
\end{equation}

Dessa forma, a partir dos dados  de fadiga coletados do algorítmo \emph{Rainflow} e de dados do material e da geometria, se tem tudo que é necessário para o desenvolvimento da vida em fadiga a partir do método de Goodman modificado. Para isso, derivam-se as expressões nos seguintes equacionamentos:

\begin{equation}
    N_f = \frac{1}{2} \left( \frac{\sigma_a}{\sigma^\prime_f - \sigma_m } \right)^{1/b}
\end{equation}

\begin{equation}
    B_f = \frac{1}{\sum \frac{1}{N_f}}
\end{equation}

Onde $B_f$ representa a vida em fadiga do material quando exposto à série  temporal de tensão utilizada na análise \emph{Rainflow}

\subsection{Método KPLS de predição de séries temporais}

O modelo KPLS (Kriging by partial least squares) é bem mais rápido que Kriging clássicos por lidar com um menor número de hiperparâmetros. Além disso ele apresenta boa performance em problemas multidimensionais. 

O equacionamento principal deste método pode ser observado na Eq.\ref{kpls}:
    
\begin{equation}\label{kpls}
    \prod^{nx}_{l=1} \prod^{nx}_{k=1} \text{exp} \left( -\theta_l(w_l^{(k)} x_l^{(i)} - w_l^{(k)} x_l^{(j)} )^2\right).
\end{equation}

O método foi aplicado a partir do pacote SMT \cite{BOUHLEL2019}.

\section{Resultados}

Neste segmento se falará dos resultados obtidos a partir dos códigos desenvolvidos em cada etapa desta pesquisa, começando com uma amostragem dos dados de tensão. Segue-se então para os resultados de vida em fadiga calculados para todos os elementos estruturais do \emph{riser}. Em seguida mostra-se o resultado de uma predição da série feita com a metodologia KPLS, com erros estatísticos associados. Mostram-se então os gráficos de vida em fadiga advindos destas séries parcialmente preditas, concluindo-se com análise dos erros.

\subsection{Estudos introdutórios}

De início aplicou-se o \emph{Rainflow} nos resultados da simulação do Anflex. Se configurou o programa para que armazenasse em memória a tensão externa de Von Mises de cada elemento estrutural para cada passo de tempo. O resultado de uma hora de simulação foram séries temporais de $18000$ passos para cada elemento estrutural. A unidade de medida das tensões era $kPa$ e o passo de tempo era de $0.2 s$. Na Fig.\ref{serie_temporal} podemos observar uma destas séries:

\begin{figure}[!ht]
    \centering
    \includegraphics[width=0.96\linewidth, trim = {0mm 0mm 0mm 0mm} , clip]{felipe/fig_felipe/tensao_tempo.png}
    \caption{Série temporal de tensões externas de Von Mises do elemento \emph{PP\_M0520212}. }
    \label{serie_temporal}
\end{figure}

A partir das séries de tensão foi possível calcular as vidas em fadiga de cada elemento com o método \emph{Rainflow} depois de se tratar os dados como descrito em \ref{rainflow}. O resultado pode ser visto adiante, junto de dados estatísticos:

\begin{figure}[!ht]
    \centering
    \includegraphics[width=\linewidth, trim = {0mm 0mm 0mm 0mm} , clip]{felipe/fig_felipe/Vida_em_fadiga_series_1.pdf} 
    \caption{Vida em fadiga para cada elemento estrutural do \emph{riser} simulado no Anflex. }
    \label{vida_em_fadiga}
\end{figure}

Observam-se bons resultados, sendo possível relacionar vários pontos no gráfico com as estruturas reais do \emph{riser} em configuração \emph{lazy-wave}.

\begin{figure}[!ht]
    \centering
    \includegraphics[width=\linewidth, trim = {0mm 0mm 0mm 0mm} , clip]{felipe/fig_felipe/Vida_em_fadiga_series_2.pdf} 
    \caption{Dados estatísticos das séries temporais de cada elemento estrutural. }
    \label{vida_em_fadiga_estatistica}
\end{figure}

Na Fig.\ref{vida_em_fadiga_estatistica} pode-se observar uma anomalia nos últimos elementos estruturais do \emph{riser}. Isso pode ser explicado pelo fato de que estes elementos fazem parte do acoplamento entre a estrutura estudada e o corpo rígido flutuante. Dessa forma, suas propriedades mecânicas são diferentes e não carregam representatividade física.

\subsection{Resultados preliminares do método KPLS}

Utilizaram-se os dados de rotação e deslocamento do ponto de ancoragem (elemento \emph{PP}) como entrada para o treino do modelo. Tem-se a predição de metade da série, isto é, foram utilizados 9000 passos para treino, prevendo-se o restante (50\%). O resultado obtido pode ser visto adiante:

\begin{figure}[!ht]
    \centering
    \includegraphics[width=1\linewidth, trim = {0mm 0mm 0mm 0mm} , clip]{felipe/fig_felipe/tensao_tempo_predicted_1.pdf} 
    \caption{Recorte da série temporal de tensões externas de Von Mises do elemento \emph{PP\_M0520212}. Com demonstração da extrapolação feita a partir do método \emph{KPLS}.}
    \label{serie_tensao_predicted_1}
\end{figure}

No caso exposto na Fig.\ref{serie_tensao_predicted_1}, observou-se uma média dos erros de $130.59 kPa$, que resulta em um desvio percentual de $0.09\%$ quando comparado ao valor médio da série de tensões. O erro por raiz da média dos quadrados das diferenças (norma L2, que denota o erro por desvio padrão) foi de $231.28 kPa$. Somente $0.16\%$ do valor médio da série temporal de tensões. O maior erro observado em toda a série foi de $1321.2 kPa$, que corresponde a somente $0.91\%$ do valor médio da série temporal.

Os dados das séries temporais de tensões expostas na Fig.\ref{serie_tensao_predicted_1} foram então analisadas com o método \emph{Rainflow} e Goodman modificado para o desenvolvimento das vidas em fadiga. Os resultados foram $5772770299$ ciclos de vida para a série de referência e $6143567360$ para a parcialmente predita, o que resultou em um erro percentual de $6.42\%$.

\subsection{Resultados gerais de vida em fadiga para séries feitas via KPLS}

Dessa forma, a análise foi feita para diferentes porcentagens de passos de tempo dedicados para treino ( $10\%$ , $20\%$ , $30\%$ , $40\%$ e $50\%$), sendo o restantes dos pontos conseguidos a partir do método KPLS. Além disso o estudo foi feito em todos os elemento estruturais do \emph{riser}.

\begin{figure}[!ht]
    \centering
    \includegraphics[width=1\linewidth, trim = {0mm 0mm 0mm 0mm} , clip]{felipe/fig_felipe/vida_fadiga_comparation.pdf} 
    \caption{Resultados de vida em fadiga calculados para cada elemento estrutural do \emph{riser}, para cada porcentagem de passos de tempo dedicados ao treino do modelo.}
    \label{comparation}
\end{figure}

Assim, notou-se que conforme se aumenta a porcentagem do número de passos de tempo dedicados ao treino observa-se uma diminuição no erro. Fez-se um estudo estatístico sobre estes resultados, com erros associados a cada porcentagem de treinamento. O resultado pode ser conferido adiante:


\begin{figure}[!ht]
    \centering
    \includegraphics[width=1\linewidth, trim = {0mm 0mm 0mm 0mm} , clip]{felipe/fig_felipe/fadiga_estatisticas.pdf} 
    \caption{Dados estatísticos dos resultados obtidos quanto a vida em fadiga. Com a média dos erros e desvio padrão. Também adimensionalizados ao serem divididos pelo valor médio de ciclos da vida em fadiga.}
    \label{estatistica_valores}
\end{figure}

Notaram-se erros que tendem a diminuir com o aumento da porcentagem de passos para treinamento. Com desvios percentuais significativos e de tendência não linear.

As séries temporais de tensão parcialmente preditas também foram analisadas e os resultados estatísticos também podem ser vistos em Fig.\ref{L1_comparation}, Fig.\ref{L2_comparation} 3 Fig.\ref{Li_comparation}. Nelas tonam-se mais regularidades quanto às manifestações da alterações de número de passos de tempo dedicados ao treino do KPLS. Também observa-se que a maior quantidade de erros encontram-se próximos à estrutura flutuante do \emph{riser} e próximos do local onde o objeto de estudo toca o leito do mar, algo que pode ser explicado pela complexidade envolvida nestes segmentos.

\begin{figure}[!ht]
    \centering
    \includegraphics[width=0.95\linewidth, trim = {0mm 0mm 0mm 0mm} , clip]{felipe/fig_felipe/L1_comparation.pdf} 
    \caption{Análise da média dos erros das séries temporais devido à aplicação do método KPLS de predição.}
    \label{L1_comparation}
\end{figure}

\begin{figure}[!ht]
    \centering
    \includegraphics[width=0.95\linewidth, trim = {0mm 0mm 0mm 0mm} , clip]{felipe/fig_felipe/L2_comparation.pdf} 
    \caption{Análise por RMS dos erros das séries temporais devido à aplicação do método KPLS de predição.}
    \label{L2_comparation}
\end{figure}

\begin{figure}[!ht]
    \centering
    \includegraphics[width=0.95\linewidth, trim = {0mm 0mm 0mm 0mm} , clip]{felipe/fig_felipe/Li_comparation.pdf} 
    \caption{Análise dos maiores erros das séries temporais devido à aplicação do método KPLS de predição.}
    \label{Li_comparation}
\end{figure}


\section{Conclusão}

Conclui-se que o método KPLS para predição de séries temporais é aplicável no desenvolvimento de vida em fadiga ao se levar em conta quanto de erro é admissível no caso estudado. Isso pode causar grande otimização no processo uma vez que reduz de forma significativa o tempo de simulação necessário para o desenvolvimento da análise sem grandes comprometimentos quanto à acurácia do método numérico. 

Vale ressaltar que a decisão de aplicação da metodologia KPLS também depende do tipo de análise que será feita, uma vez que a predição causa desvios significativos dependendo do segmento do \emph{riser}, como nos entornos de zonas flutuantes e nos pontos de contato das estruturas com o leito marinho. Assim, pontos sujeitos a grandes não linearidades dinâmicas resultam em anomalias perceptíveis. 

No futuro deste estudo pretende-se analisar mais porcentagens de traino para o método KPLS, além de se procurar outros casos para se aplicar esta análise.

%	\include{helio/}
%	\include{marcelo/}
%	\include{tatiane/}

	\chapter{Conclusões e considerações finais}
		\label{finalsec}
		
		Este relatório técnico parcial teve como objetivo apresentar os desenvolvimentos realizados no período de junho de 2019 -- dezembro de 2019. Entre as principais contribuições compreendidas no mesmo, pode-se citar a revisão bibliográfica acerca de vários métodos de metamodelagem, a implementação computacional de algumas destas técnicas, bem como de procedimentos de amostragem, com posterior aplicação das mesmas em funções de teste bem difundidas na literatura, o êxito na realização de prolongamentos de séries temporais advindas de simulações realizadas com o \textit{software} Anflex e, por fim, a validação da ferramenta computacional MFSim através da comparação dos sinais de CD e CL obtidos com dados provenientes de experimentação material.
		
		\begin{equation}
			\frac{1}{2}
		\end{equation}







		
	\begin{singlespace}
		\phantomsection
		\addcontentsline{toc}{chapter}{Referências Bibliográficas}
		 %\bibliographystyle{acm}
		% \bibliographystyle{plain}
		\bibliography{metamodelo}
	\end{singlespace}
	%\include{apendices}

\end{document}   % Fim.
